% Part: conditional-logics
% Chapter: minimal-change-semantics
% Section: true-false

\documentclass[../../../include/open-logic-section]{subfiles}

\begin{document}

\olfileid{con}{min}{tf}

\olsection{صدق و کذب شرطی‌های خلاف واقع}

شرطی خلاف واقع~$!A \cif !B$ هنگامی (به‌طور غیرتهی) صادق است که همهٔ
نزدیک‌ترین جهان‌های~$!A$، جهان‌های~$!B$ باشند، چنان‌که در
\olref{fig:true} نشان داده شده است.
\begin{figure}
\begin{center}
\begin{tikzpicture}[scale=.7]\small
  \spheresystem{5}
  \draw (0,0) node {$w$};
  \propositionintersect{0}{4}{40}{3.5}
  {\tikzset{proposition/.append={smooth,tension=1.4}}
  \proposition[shift={(1.6,-1)}]{90}{1}{30}{4}}
  \path (1.5,3) node[below] {$\formula{B}$};
  \path (3,0) node {$\formula{A}$};
\end{tikzpicture}
\caption{شرطی خلاف واقعِ به‌طور غیرتهی صادق}
\ollabel{fig:true}
\end{center}
\end{figure}
شرطی خلاف واقع در~$w$ همچنین هنگامی صادق است که دستگاه کره‌های پیرامون
$w$ اصلاً هیچ کرهٔ پذیرای~$!A$ نداشته باشد. در این حالت شرطی
\emph{به‌طور تهی} صادق است (بنگرید به~\olref{fig:vacuous}).
\begin{figure}
\begin{center}
\begin{tikzpicture}[scale=.7]\small
  \spheresystem{5}
  \draw (0,0) node {$w$};
  \proposition{0}{6.5}{40}{3.5}
  {\tikzset{proposition/.append={smooth,tension=1.4}}
  \proposition[shift={(1.2,-1)}]{90}{1}{30}{4}}
  \path (1.5,3) node[below] {$\formula{B}$};
  \spherepos{0}{7}{node {$\formula{A}$}}
\end{tikzpicture}
\caption{شرطی خلاف واقعِ به‌طور تهی صادق}
\ollabel{fig:vacuous}
\end{center}
\end{figure}

شرطی خلاف واقع به دو شیوه می‌تواند کاذب باشد. یک شیوه آن است که برخی،
اما نه همهٔ، نزدیک‌ترین جهان‌های~$!A$ جهان‌های~$!B$ باشند. در این حالت
$!A \cif \lnot !B$ نیز کاذب است (بنگرید به~\olref{fig:false}).
\begin{figure}
\begin{center}
\begin{tikzpicture}[scale=.7]\small
  \spheresystem{5}
  \draw (0,0) node {$w$};
  \propositionintersect{0}{4}{40}{3.5}
  {\tikzset{proposition/.append={smooth,tension=1.4}}
  \proposition[shift={(1.6,0)}]{90}{1}{30}{3}}
  \path (1.5,3) node[below] {$\formula{B}$};
  \path (3,0) node {$\formula{A}$};
\end{tikzpicture}
\caption{شرطی خلاف واقع کاذب، مقابل کاذب}
\ollabel{fig:false}
\end{center}
\end{figure}
اگر نزدیک‌ترین جهان‌های~$!A$ اصلاً با جهان‌های~$!B$ هم‌پوشانی نداشته
باشند، آنگاه~$!A \cif !B$ کاذب است. اما در این حالت همهٔ نزدیک‌ترین
جهان‌های~$!A$، جهان‌های~$\lnot !B$ هستند و بنابراین
$!A \cif \lnot !B$ صادق است (بنگرید به~\olref{fig:false-opposite}).
\begin{figure}
\begin{center}
\begin{tikzpicture}[scale=.7]\small
  \spheresystem{5}
  \draw (0,0) node {$w$};
  \propositionintersect{0}{4}{40}{3.5}
  {\tikzset{proposition/.append={smooth,tension=1.4}}
  \proposition[shift={(-1,0)}]{90}{1}{30}{3}}
  \path (-1,3) node[below] {$\formula{B}$};
  \path (3,0) node {$\formula{A}$};
\end{tikzpicture}
\caption{شرطی خلاف واقع کاذب، مقابل صادق}
\ollabel{fig:false-opposite}
\end{center}
\end{figure}

برخلاف شرطی اکید، شرطی‌های خلاف واقع می‌توانند امکانی باشند. مدل
کره‌ایِ~\olref{fig:contingent} را در نظر بگیرید. همهٔ جهان‌های~$!A$ که
به~$u$ نزدیک‌ترین‌اند، جهان‌های~$!B$ هستند، پس
$\mSat{M}{!A \cif !B}[u]$. اما برخی جهان‌های~$!A$ که به~$v$
نزدیک‌ترین‌اند جهان‌های~$!B$ نیستند، پس~$\mSat/{M}{!A \cif !B}[v]$.
\begin{figure}
\begin{center}
\begin{tikzpicture}[scale=.6]\small
  \spheresystem{7}
  \spheresystem[shift={(0:3.1)},dashed]{4}
  \propositionintersect{45}{4}{40}{4.5}
  \begin{scope}
    \clip \propositionplot{45}{4}{40}{4.5} ;
    \spherelayer[shift={(0:3.1)}]{3}
  \end{scope}
  \proposition[shift={(1.4,-.5)}]{90}{1}{35}{5}
  \path (0,0) node {$u$};
  \path (0:3.1) node {$v$};
  \spherepos{35}{9}{node {$\formula{A}$}}
  \spherepos{80}{9}{node {$\formula{B}$}}
\end{tikzpicture}
\end{center}
\caption{شرطی خلاف واقع امکانی}
\ollabel{fig:contingent}
\end{figure}
\end{document}
